\documentclass[a4paper]{article}
\usepackage[pdftex]{graphicx}
\usepackage[utf8]{inputenc}
\usepackage{enumerate}
\usepackage{siunitx}
\sisetup{locale=DE}
\usepackage{href-ul}
\hypersetup{
	colorlinks=true,
	linkcolor=blue,
	urlcolor=blue}
\usepackage{icomma}
\usepackage{amssymb}
\usepackage{geometry}
\geometry{a4paper, top=15mm, left=15mm, right=15mm, bottom=15mm,
	headsep=10mm, footskip=12mm}


\begin{document}
	%Title
	{\bf Physics school assignment}\\
	\begin{enumerate}[1.]
		
		{\bf \item Laser light of wavelength $\lambda=\SI{546}{\nano\meter}$ }falls on a double slit. Bright stripes are observed on a screen 2.80 m away.
			\begin{enumerate}[a)]
				\item On the screen, the two 1st order maxima are 17.6 cm apart.\\
				Calculate the gap distance.
				\item How do the number and position of the maxima change if laser light of a smaller wavelength is used with an otherwise unchanged arrangement?
			\end{enumerate}
			
			{\bf \item How long would the screen distance} have to be at least to reach the sodium double line ($\SI{588.99}{\nano\meter}$ and $\SI{589.59}{\nano\meter} $) with a grid of the slit distances $b=\SI{100}{\micro\meter}$, so that the distance between the two first-order maxima is $\SI{3}{\milli\meter}$?
			
			
		\end{enumerate}
		
		\fbox{
			\begin{minipage}{0.5\textwidth}
				For the solution, please \href{https://www.okuyakl.de/phys/p11wellL095/le095.pdf}{click here} or scan the QR code.\\
				Additional worksheets are available at
				
				\href{http://www.okuyakl.com}{www.okuyakl.com}
			\end{minipage}
			\hfill
			\begin{minipage}{0.4\textwidth}
				\includegraphics[width=1.5 cm]{../../viecher/zwe03}
				\includegraphics[width=3 cm]{qre095}
				\includegraphics[width=2 cm]{../../viecher/afanticon1}
				
		\end{minipage}}
		
	\end{document}